% !TeX root = ../../Book3.tex
\chapter{Green 函数、Green 域与多连通区域}
\label{b3:ch:17}
\BookChapterNumber{b3:ch:17}
% 本章摘要（不计入编号）
\begin{chapterabstract}
Green 函数把对数点源与零边界约束合在一个正核中。本章由 Perron 解和穷竭构造核，证明最小性、对称性、共形不变性和区域单调性；随后讨论 Green 势、容量与多连通域的周期分解。
\end{chapterabstract}

% 本章目标（不计入编号）
\begin{learninggoals}
能核对圆盘和半平面 Green 核的符号与常数；能利用穷竭和正核上界判断存在性；能区分 Green 域与覆盖意义的双曲型；能用边界调和测度调整有限连通域的共轭周期。
\end{learninggoals}

圆盘中心的函数 \(\log(R/|z|)\) 在边界为零，在原点有一个正无穷点源。把半径增大，内部势就增加；若一直增大到全平面，这个核发散，说明域的大小会影响点源能否有有限的零边界响应。

\ProseParagraphBreak

本章以\chapref{b3:ch:15}的调和比较和\chapref{b3:ch:16}的 Perron 方法为前提。一般势论与容量的进一步阅读为 Ransford 的《Potential Theory in the Complex Plane》\cite{Ransford1995Potential}与 Armitage、Gardiner 的《Classical Potential Theory》\cite{Armitage2001Classical}；仅核书目和主题的部分不承担这里的证明来源。

% !TeX root = ../../Book3.tex
\section{Green 函数}
\label{b3:sec:1701}

在区域内放置一个点源，并让边界保持零势，需要同时安排局部奇性和全边界条件。对数基本解给出点源，Dirichlet 问题则补上一个调和修正。本章采用正 Green 核，奇性为 \(-\log|z-w|\)，因而其 Laplace 算子带负号。先在有界正则域完成这项构造，再说明光滑边界如何用于核的对称性。

\subsection{定义}
\label{b3:subsec:170101}

设 \(D\subseteq\mathbb C\) 为有界区域，\(w\in D\)。若实函数 \(g_D(\cdot,w)\) 在 \(D\setminus\{w\}\) 严格为正且调和，在每个 \(\xi\in\partial D\) 满足
\[
 \lim_{D\ni z\to\xi}g_D(z,w)=0,
\]
并且 \(g_D(z,w)+\log|z-w|\) 可在 \(w\) 附近延拓为调和函数，则称它为以 \(w\) 为极点的经典零边值\term[green-hanshu]{Green 函数}[Green function]。局部奇性中的系数固定为一；只规定趋于正无穷，不能确定点源的大小。

\ProseParagraphBreak
记调和余项为 \(h_w\)，则在整个 \(D\) 内有
\[
 g_D(z,w)=-\log|z-w|+h_w(z),\quad z\ne w.
\]
这里 \(h_w\) 在极点外由等式定义，在极点处由假设延拓；因此它在全域调和。由\cref{b3:prop:15-log-fundamental}，对数函数局部可积且
\(\Delta\log|z-w|=2\pi\delta_w\)，所以
\[
 -\Delta_zg_D(\cdot,w)=2\pi\delta_w
\]
按分布成立。这个等式记录点源，零边值则记录区域的约束，两项都属于核的定义。

\ProseParagraphBreak
为了在边界使两项相消，所求修正必须满足
\(h_w(\xi)=\log|\xi-w|\)。因此使用的是正号的对数边值。因为 \(w\) 在内部，它到边界有正距离，这份边值在紧集 \(\partial D\) 上连续且有界；上一章的连续 Dirichlet 理论正好可以处理它。

\Needspace{6\baselineskip}
\begin{theorem}[有界正则域上的 Green 核]\label{b3:thm:17-green-regular}
设 \(D\subseteq\mathbb C\) 为有界正则域，即每个连续实边值的 Perron 解都连续达到全部边界值。对 \(w\in D\)，令
\[
 \varphi_w(\xi)=\log|\xi-w|,\quad \xi\in\partial D,
 \quad h_w=H_{\varphi_w}.
\]
则
\[
 g_D(z,w)=-\log|z-w|+h_w(z),\quad z\in D\setminus\{w\},
\]
是以 \(w\) 为极点的经典零边值 Green 函数。
\end{theorem}
\begin{proof}
极点到边界的正距离保证 \(\varphi_w\in C(\partial D;\mathbb R)\)。
\cref{b3:thm:16-perron}给出内部调和的 Perron 解，而 \(D\) 的正则性保证
\(h_w\in C(\overline D)\) 且 \(h_w|_{\partial D}=\varphi_w\)。因此所定义的
\(g_D\) 在极点外调和，连续取零边值，其对数奇性也符合定义。

还须证明严格正性。先取 \(\rho>0\) 使
\(\overline{B(w,\rho)}\subset D\)。调和函数 \(h_w\) 在这个闭圆盘上有界，故
\[
 \min_{|z-w|=\varepsilon}g_D(z,w)
 \ge-\log\varepsilon-\sup_{\overline{B(w,\rho)}}|h_w|
 \longrightarrow+\infty
 \quad(\varepsilon\downarrow0).
\]
取足够小的 \(\varepsilon\)，在
\(D\setminus\overline{B(w,\varepsilon)}\) 上对 \(-g_D\) 应用边界最大值原理：原边界上它为零，小圆上它为负，所以 \(g_D\ge0\) 于余域；若余域有多个连通分支，就在每个分支分别使用同一原理。任意固定 \(z\ne w\) 都可放进这样的余域，故全域非负。

区域 \(D\setminus\{w\}\) 仍连通：连接两点的内部折线若经过 \(w\)，可在一个内部小圆盘中以圆弧绕开该点。若 \(g_D\) 在某个非极点位置为零，强最小值原理便使它在整个 \(D\setminus\{w\}\) 恒零，与极点附近趋于正无穷矛盾。因此 \(g_D>0\)。
\end{proof}

构造只用正则性，没有用边界导数。若记
\[
 d_w=\min_{\xi\in\partial D}|\xi-w|,\quad
 L_w=\max_{\xi\in\partial D}|\xi-w|,
\]
则 \(0<d_w\le L_w<\infty\)，最大值原理给出
\[
 \log d_w\le h_w(z)\le\log L_w,\quad z\in D.
\]
因此极点附近有明确的比较
\[
 \log\frac{d_w}{|z-w|}
 \le g_D(z,w)\le\log\frac{L_w}{|z-w|}.
\]
左端离极点较远时可以为负，严格正性仍由上面的域内比较证明。在圆盘
\(B(w,R)\) 中，边值 \(\varphi_w\) 恒等于 \(\log R\)，所以
\(h_w=\log R\)，得到 \(g_{B(w,R)}(z,w)=\log(R/|z-w|)\)。

\subsection{唯一性}
\label{b3:subsec:170102}

同一对数奇性的差不再含点源，因此经典核的唯一性仍归结为调和函数的零边值唯一性。这一步只要求区域有界；正则性是保证存在性的条件。

\Needspace{6\baselineskip}
\begin{proposition}[经典零边值核的唯一性]\label{b3:prop:17-green-uniqueness}
设 \(D\subseteq\mathbb C\) 为有界区域，\(w\in D\)。在
\(D\setminus\{w\}\) 调和、在每个边界点连续取零值，且
\(g(z)+\log|z-w|\) 在 \(w\) 附近调和可去的实函数 \(g\) 至多一个。
这里不必预先假定候选函数为正。
\end{proposition}
\begin{proof}
若 \(g_1,g_2\) 是两个候选，则在极点附近写成
\(g_j=-\log|z-w|+h_j\)。它们的差等于 \(h_1-h_2\)，故可在 \(w\) 处延拓为调和函数。延拓后的差在全域调和，并在每个边界点趋于零。由有界域的边界上极限原理\cref{b3:cor:15-boundary-limsup}，对差与负差分别比较，得到 \(g_1-g_2=0\)。
\end{proof}

经典核还有一个不要求竞争者满足零边值的性质。它将用于把定义扩展到边界粗糙或无界的区域。

\Needspace{6\baselineskip}
\begin{proposition}[经典 Green 核的最小性]\label{b3:prop:17-classical-green-minimal}
设 \(D\subseteq\mathbb C\) 为有界正则域，\(w\in D\)，\(g_D(\cdot,w)\) 为其经典 Green 核。若实函数 \(G\) 在 \(D\setminus\{w\}\) 为正且调和，并且 \(G(z)+\log|z-w|\) 在 \(w\) 附近调和可去，则
\[
 g_D(z,w)\le G(z),\quad z\in D\setminus\{w\}.
\]
对 \(G\) 不要求边界极限存在，也不要求它在靠近边界时有界。
\end{proposition}
\begin{proof}
函数 \(q=g_D-G\) 的奇性相消，故延拓为 \(D\) 上的调和函数。因为
\(G>0\)，对每个 \(\xi\in\partial D\) 都有
\[
 \limsup_{D\ni z\to\xi}q(z)
 \le\lim_{D\ni z\to\xi}g_D(z,w)=0.
\]
区域有界，\cref{b3:cor:15-boundary-limsup}直接给 \(q\le0\)，即得结论。
\end{proof}

最小性比较的是同一对数奇性的全部正竞争函数。若向经典核加一个正的调和函数，奇性与正性仍然保留，但一般不再取零边值，也不再是最小者。对无界域，甚至通常边界上的零值也可能保留：上半平面的函数
\(y\) 为正、调和且在实轴为零，可以加到任何已有的正零边值候选上。因此后面的穷竭构造将以最小性识别一般 Green 核，而不要求它在每个不正则边界点取零。

\ProseParagraphBreak
以下边界积分限定在有限条实解析 Jordan 曲线围成的有界域。这里曲线的实解析参数要求导数处处非零，因而排除了尖点。这类边界局部是 \(C^1\) 图像，域的补集包含一个外锥，由\cref{b3:prop:16-exterior-cone-regular}可知域正则。零边值的反射还会给出法向导数，不能在证明反射之前先假定这些导数存在。

\Needspace{6\baselineskip}
\begin{lemma}[常边值调和函数的局部反射]\label{b3:lem:17-harmonic-reflection}
设 \(D\subseteq\mathbb C\) 为有界区域，其边界由有限条互不相交的正则实解析 Jordan 曲线组成。设 \(I\) 是一段开边界弧，\(V\subseteq\mathbb C\) 是包含 \(I\) 的开集，实函数 \(u\) 在 \(D\cap V\) 调和，并在 \(I\) 上连续取同一常值 \(c\)。则对每个 \(\xi\in I\)，\(u\) 都可调和延拓到包含 \(\xi\) 的一个开邻域。因此 \(u\) 在该弧局部光滑直到边界。
\end{lemma}
\begin{proof}
在 \(\xi\) 附近取实解析参数 \(\gamma(t)\)，使 \(\gamma(0)=\xi\)、
\(\gamma'(0)\ne0\)。其幂级数在一个复圆盘中给出全纯延拓
\(\Gamma(\zeta)\)。由复逆函数定理，缩小圆盘后 \(\Gamma\) 双全纯，且将一段实轴映为边界弧。正则 Jordan 边界在局部把邻域分成两侧；调整参数方向后，可使上半圆盘映到 \(D\) 侧。选 \(a>0\) 足够小，使闭上半圆盘
\(\{|\zeta|\le a,\operatorname{Im}\zeta\ge0\}\) 都落在这张坐标的范围内，且半圆弧除两端外映入 \(D\)。

令 \(v(\zeta)=u(\Gamma(\zeta))-c\)。链式法则给
\[
 \Delta v=|\Gamma'|^2\cdot(\Delta u)\circ\Gamma=0
\]
于上半圆盘。它连续到闭上半圆盘，并在直径上为零。在完整圆周
\(|\zeta|=a\) 上定义数据
\[
 \psi(\zeta)=
 \begin{cases}
 v(\zeta),&\operatorname{Im}\zeta\ge0,\\
 -v(\overline\zeta),&\operatorname{Im}\zeta<0.
 \end{cases}
\]
两端 \(\zeta=\pm a\) 的值为零，因此 \(\psi\) 是连续圆周数据。由
\cref{b3:thm:16-poisson}，其 Poisson 解 \(V\) 在整圆盘调和并连续到圆周。

函数 \(-V(\overline\zeta)\) 具有相同边值 \(\psi\)，圆盘解的唯一性给
\(V(\zeta)=-V(\overline\zeta)\)，所以 \(V\) 在直径上为零。在上半圆盘，
\(V\) 与 \(v\) 都调和，且在半圆弧与直径上取相同值；有界域最大值原理给
\(V=v\)。于是
\[
 c+V\circ\Gamma^{-1}
\]
在 \(\xi\) 的完整邻域中调和，并在域内与原来的 \(u\) 相同。这就是所需延拓。
\end{proof}

对固定极点 \(w\)，边界小邻域可选得不含 \(w\)，故该引理适用于
\(g_D(\cdot,w)\)。边界紧且分量有限，取有限个这样的邻域就得到整个边界上的光滑性及法向导数的连续性。它们足以在删去极点小圆后使用 Green 公式；一般粗糙域的对称性将由穷竭传递，不在粗糙边界上写通常法向积分。

\subsection{对称性}
\label{b3:subsec:170103}

Green 核描述源点与观察点之间的作用。交换二者以后核值不变，需要把两份核放进同一个积分恒等式；两处奇性分别贡献一个函数值。

\Needspace{6\baselineskip}
\begin{proposition}[Green 核的对称性]\label{b3:prop:17-green-symmetry}
设 \(D\subseteq\mathbb C\) 为有界区域，其边界由有限条互不相交的正则实解析 Jordan 曲线组成，\(g_D\) 为对数奇性取 \(-\log|z-w|\) 的经典零边值 Green 核。则对不同的 \(z,w\in D\)，有
\[
 g_D(z,w)=g_D(w,z).
\]
\end{proposition}
\begin{proof}
固定 \(z\ne w\)，在变量 \(\zeta\) 中记
\(u(\zeta)=g_D(\zeta,z)\)、\(v(\zeta)=g_D(\zeta,w)\)。取足够小的
\(\varepsilon>0\)，使以 \(z,w\) 为中心的两个闭圆盘互不相交且均包含于
\(D\)。令
\[
 \Omega_\varepsilon
 =D\setminus\bigl(\overline{B(z,\varepsilon)}
                  \cup\overline{B(w,\varepsilon)}\bigr).
\]
两函数在这个域中调和，且由边界反射光滑到原边界。这里使用第二卷推论1.2.3的 Green 公式：区域有界、边界为 \(C^1\)，函数在闭域的邻域中为 \(C^2\)。原边界实解析，新增小圆也光滑；反射与极点外的光滑性给出所需邻域延拓，故这些条件均满足。对恒等式
\[
 \operatorname{div}(u\nabla v-v\nabla u)=u\Delta v-v\Delta u
\]
使用平面散度公式，得到 Green 第二公式
\[
 0=\int_{\partial\Omega_\varepsilon}
       (u\,\partial_nv-v\,\partial_nu)\,\mathrm ds,
\]
其中 \(n\) 一律指 \(\Omega_\varepsilon\) 的外单位法向。原边界的每个分量上
\(u=v=0\)，故那里没有贡献。

在 \(z\) 附近写 \(u(\zeta)=-\log|\zeta-z|+a(\zeta)\)，其中 \(a\) 光滑调和；函数 \(v\) 在这里也光滑。小圆上的外法向朝向被删圆盘的中心，故
\[
 \begin{aligned}
 u&=-\log\varepsilon+O(1),&
 \partial_nu&=\frac1\varepsilon+O(1),\\
 v&=v(z)+O(\varepsilon),&
 \partial_nv&=O(1).
 \end{aligned}
\]
所有误差均对圆周方向一致。因此
\[
 \int_{\partial B(z,\varepsilon)}
       (u\,\partial_nv-v\,\partial_nu)\,\mathrm ds
 =-2\pi v(z)+O(\varepsilon|\log\varepsilon|)+O(\varepsilon).
\]
第一项 \(u\,\partial_nv\) 乘上圆周长后趋于零；保留下来的是带负号的
\(-v\,\partial_nu\)。

在 \(w\) 的小圆上两函数的角色交换：
\(v=-\log\varepsilon+O(1)\)、\(\partial_nv=1/\varepsilon+O(1)\)，而
\(u=u(w)+O(\varepsilon)\)、\(\partial_nu=O(1)\)。这里得到
\[
 \int_{\partial B(w,\varepsilon)}
       (u\,\partial_nv-v\,\partial_nu)\,\mathrm ds
 =2\pi u(w)+O(\varepsilon|\log\varepsilon|)+O(\varepsilon).
\]
两小圆积分之和为零，令 \(\varepsilon\downarrow0\)，得
\(u(w)=v(z)\)，即所述对称性。
\end{proof}

小圆上的法向不是从其中心向外，而是余域的向外方向，即朝向被删掉的圆盘。这个方向使对数奇项的法向导数为 \(+1/\varepsilon\)；两处主项的相反符号则来自
\(u\,\partial_nv-v\,\partial_nu\) 中两函数的位置。对称性还表明，对固定
\(z\)，函数 \(w\mapsto g_D(z,w)\) 在 \(D\setminus\{z\}\) 调和，因为它等于已按第一个变量构造的 \(g_D(w,z)\)。

\subsection{共形不变性}
\label{b3:subsec:170104}

共形变换保持二维调和方程，但对数奇性会多出一个调和项。经典核的边界条件还需要映射在闭包上的对应；先保留这项条件，可以直接使用唯一性。

\Needspace{6\baselineskip}
\begin{proposition}[经典 Green 核的共形不变性]\label{b3:prop:17-classical-green-conformal}
设 \(D,D'\subseteq\mathbb C\) 为有界正则域，\(F:D\to D'\) 双全纯，且
\(F\) 延拓为 \(\overline D\) 到 \(\overline{D'}\) 的同胚。则对不同的
\(z,w\in D\)，两域的经典 Green 核满足
\[
 g_D(z,w)=g_{D'}(F(z),F(w)).
\]
\end{proposition}
\begin{proof}
固定 \(w\)，令 \(G(z)=g_{D'}(F(z),F(w))\)。它为正，并由调和函数复合公式在
\(D\setminus\{w\}\) 调和。闭包同胚将 \(\partial D\) 对应到
\(\partial D'\)，故 \(G\) 在每个边界点趋于零。

在 \(w\) 的小圆盘上写
\[
 F(z)-F(w)=(z-w)q(z),\quad q(w)=F'(w)\ne0.
\]
缩小圆盘使全纯函数 \(q\) 处处非零，于是 \(\log|q|\) 调和。若
\(g_{D'}(\eta,F(w))=-\log|\eta-F(w)|+h(\eta)\)，则
\[
 G(z)+\log|z-w|=h(F(z))-\log|q(z)|
\]
在 \(w\) 附近调和。因而 \(G\) 具有正确的奇性和零边值，由
\cref{b3:prop:17-green-uniqueness}得到 \(G=g_D(\cdot,w)\)。
\end{proof}

映射在边界上不必有非零导数；本命题只在那里使用同胚和连续性，导数非零只用在内部极点。相似变换也说明为何核不能只看奇性中的距离：若
\(F(z)=az+b\)、\(a\ne0\)，则
\(-\log|F(z)-F(w)|=-\log|z-w|-\log|a|\)，调和修正必须补回这项常数，才能保持整个核值不变。

\ProseParagraphBreak
没有通常边界延拓的内部共形双射不满足本命题的边界假设。一般 Green 域仍有同一不变性，但需先建立最小正核，再把正竞争函数双向拉回进行比较；这条不使用边界延拓的证明将在\cref{b3:sec:1703}中给出。

% !TeX root = ../../Book3.tex
\section{Green 域与抛物型域}
\label{b3:sec:1702}

上一节在有界正则域中把点源与连续零边值配在一起。若区域无界，或者边界含有不正则点，就不能继续把通常的边界值当作唯一的约束。本节先问一个更基本的问题：在整个区域内，能否有一个始终为正、只有指定对数奇性的函数？下一节再从所有这样的函数中构造最小者。

\subsection{Green 域}
\label{b3:subsec:170201}

为了在尚未讨论无穷远行为时准确表述存在性，先把竞争函数类写出来。

\Needspace{6\baselineskip}
\begin{definition}[正基本核与 Green 域]\label{b3:def:17-green-domain}
设 \(D\subseteq\mathbb C\) 为区域，\(w\in D\)。称实函数 \(q:D\setminus\{w\}\to(0,\infty)\) 为极点 \(w\) 的正基本核，如果它在 \(D\setminus\{w\}\) 调和，且 \(q(z)+\log|z-w|\) 在 \(w\) 的某个圆盘邻域中延拓为调和函数。记所有正基本核的集合为 \(\mathcal Q_D(w)\)。若对每个 \(w\in D\) 都有 \(\mathcal Q_D(w)\ne\varnothing\)，称 \(D\) 为\term[green-yu]{Green 域}[Greenian domain]。当 \(\mathcal Q_D(w)\) 中存在逐点不超过其余各项的成员时，称这个最小成员为一般域上的 Green 核，记为 \(g_D(\cdot,w)\)。
\end{definition}

“最小”指与全体竞争函数逐点比较，不是只在某个内点使数值最小。一个类如果有最小成员，该成员必唯一。上一节已经证明，有界正则域的经典 Green 核就是其正基本核类的最小成员。下一节将证明，只要该类非空，最小成员就存在；还将证明一处极点的存在性足以推出全部极点的存在性。因此这里的定义与用有限穷竭极限定义 Green 域等价，但现在不把尚未构造的极限作为论证前提。

\ProseParagraphBreak
有界域不必正则，也一定有正基本核。固定 \(w\in D\)，选择 \(R>0\) 使 \(D\subset B(w,R)\)，则
\[
 q(z)=\log\frac R{|z-w|}
\]
在 \(D\setminus\{w\}\) 严格为正、调和，且 \(q+\log|z-w|=\log R\)。所以每个有界区域都是 Green 域。这个竞争核一般不在 \(\partial D\) 上取零；其用途是给随后构造的最小核一个有限上界。边界是否正则影响连续零边值的达到，不决定有界域是否有 Green 核。

\ProseParagraphBreak
正基本核也可以限制到子域。若 \(D\subset E\)，\(w\in D\)，则 \(\mathcal Q_E(w)\) 的每一项在 \(D\) 上仍有正确奇性和正性。因此 Green 域的每个子域也是 Green 域。这一判断只用竞争函数的存在，尚不需要最小核的区域单调性。

\subsection{Green 函数不存在的情形}
\label{b3:subsec:170202}

全平面的困难出现在无穷远。假设极点 \(w\) 有一个正基本核 \(q\)，则
\[
 h(z)=q(z)+\log|z-w|
\]
延拓为整个平面的调和函数。在以 \(w\) 为中心的圆周上取平均，得到
\[
 \frac1{2\pi}\int_0^{2\pi}q(w+re^{\mathrm it})\,\mathrm dt
 =h(w)-\log r.
\]
当 \(r\) 足够大时右端为负，与 \(q>0\) 矛盾。因此整个复平面没有正基本核，更没有 Green 核。这里不是因为没有可写的基本解：\(-\log|z-w|\) 正是所需基本解，只是任何整平面调和修正都不能使它在全部半径上保持正平均。

\ProseParagraphBreak
删除有限个点是否可以补救这一困难，需要先控制那些点附近可能出现的奇性。对于任意调和函数，孤立奇性可以比对数强得多；正性却排除了这些振荡主部。

\Needspace{6\baselineskip}
\begin{lemma}[Bôcher 孤立奇性分解]\label{b3:lem:17-positive-singularity}
设 \(R>0\)，\(u\ge0\) 在 \(0<|z|<R\) 调和。则存在唯一的 \(c\ge0\) 及整个圆盘 \(B(0,R)\) 上的实调和函数 \(h\)，使
\[
 u(z)=-c\log|z|+h(z)\quad(0<|z|<R).
\]
\end{lemma}
\begin{proof}
对每个整数 \(n\) 定义圆周系数
\[
 a_n(r)=\frac1{2\pi}\int_0^{2\pi}u(re^{\mathrm it})e^{-\mathrm int}\,\mathrm dt,
 \quad 0<r<R.
\]
任意闭圆环内的 \(u\) 光滑，故可以对 \(r\) 求导；将极坐标 Laplace 方程乘以 \(e^{-\mathrm int}\)，并对角变量积分两次分部，得到
\[
 a_n''(r)+\frac1r a_n'(r)-\frac{n^2}{r^2}a_n(r)=0.
\]
当 \(n=0\) 时，解为 \(a_0(r)=A+B\log r\)。圆周平均非负，令 \(r\downarrow0\) 可知 \(B\le0\)。当 \(n\ne0\) 时，换元 \(x=\log r\) 将方程化为常系数方程，因而
\[
 a_n(r)=\alpha_n r^{|n|}+\beta_n r^{-|n|}.
\]
正性给出关键估计
\[
 |a_n(r)|\le\frac1{2\pi}\int_0^{2\pi}u(re^{\mathrm it})\,\mathrm dt
 =a_0(r)=O(1+|\log r|).
\]
两端乘以 \(r^{|n|}\)，再令 \(r\downarrow0\)，便知 \(\beta_n=0\)。负幂的增长比对数快，因而不能在正性所给的圆周积分界下保留。

固定 \(0<S<R\)，令 \(h_S\) 为圆周数据 \(u(Se^{\mathrm it})-B\log S\) 的 Poisson 延拓。零频率系数为 \(A\)，非零系数为 \(\alpha_n S^{|n|}\)。由\cref{b3:prop:16-poisson-frequencies}，在任意 \(0<r<S\) 上，\(h_S(re^{\mathrm it})\) 的系数恰为 \(A\) 及 \(\alpha_n r^{|n|}\)。这些也正是 \(u(re^{\mathrm it})-B\log r\) 的全部系数。

两个连续圆周函数的所有 Fourier 系数相同，就必相等。为提醒这一唯一性的依据，第16章圆环构造已证明连续周期函数的 Fej\'er 平均一致收敛；如果差函数的全部系数为零，它的每个 Fej\'er 平均都为零，因此差函数为零。于是 \(u-B\log|z|=h_S\) 于 \(0<|z|<S\)。不同 \(S\) 的调和延拓在交叠的去心圆盘一致，连续性使原点值也相同，故拼成整个 \(B(0,R)\) 上的调和函数 \(h\)。取 \(c=-B\ge0\) 得到表示。零频率确定 \(B\)，其余函数由差确定，所以表示唯一。
\end{proof}

这个结论也称正调和函数的\term[bocher-gulixing-fenjie]{Bôcher 孤立奇性分解}[Bôcher decomposition of an isolated singularity]，这里证明的是二维版本。若允许 \(u\) 改变符号，\(\operatorname{Re}(1/z)\) 就有负幂奇性；它的圆周平均为零，绝对值积分却很大，证明中的 \(|a_n|\le a_0\) 不再成立。

\Needspace{6\baselineskip}
\begin{proposition}[有限穿孔平面无正 Green 核]\label{b3:prop:17-finite-punctures-parabolic}
设 \(A\subset\mathbb C\) 为有限集，允许为空。则在 \(D=\mathbb C\setminus A\) 上，对任何 \(w\in D\) 都不存在正基本核。
\end{proposition}
\begin{proof}
空集情形已由圆周平均证明。设 \(A=\{a_1,\ldots,a_m\}\)，假设 \(q\in\mathcal Q_D(w)\)。在每个 \(a_j\) 的一个不含其余缺口及 \(w\) 的小圆盘内，上一引理给
\[
 q(z)=-c_j\log|z-a_j|+h_j(z),\quad c_j\ge0,
\]
其中 \(h_j\) 跨过 \(a_j\) 调和。因此
\[
 H(z)=q(z)+\log|z-w|+\sum_{j=1}^m c_j\log|z-a_j|
\]
延拓为整平面调和函数：在每个被删点，恰由相应局部分解消去奇性；在 \(w\) 则由正基本核的定义消去奇性。

取 \(r\) 大于 \(|w|,|a_1|,\ldots,|a_m|\)。对任意 \(|a|<r\)，有
\[
 \frac1{2\pi}\int_0^{2\pi}\log|re^{\mathrm it}-a|\,\mathrm dt=\log r.
\]
这可由 \(\log|re^{\mathrm it}-a|=\log r+\operatorname{Re}\log(1-ae^{-\mathrm it}/r)\) 的一致收敛幂级数直接积分，所有非零频率平均为零。于是
\[
 \frac1{2\pi}\int_0^{2\pi}q(re^{\mathrm it})\,\mathrm dt
 =H(0)-\left(1+\sum_{j=1}^m c_j\right)\log r.
\]
右端最终为负，左端始终为正，矛盾。
\end{proof}

每个缺口允许新增的系数均为非负，它们只能使大圆平均中的负对数斜率更大。因而不能靠有限个孤立缺口补偿无穷远的正性失败。这也是论证中必须分析缺口奇性，而不能只说“孤立点没有影响”的原因。

\subsection{\texorpdfstring{\(\mathbb C\)}{C} 的抛物性}
\label{b3:subsec:170203}

从势论角度，在任何极点都没有正基本核的区域称为\term[paowuxing-yu]{抛物型域}[parabolic domain]。这里的名称针对 Green 核的不存在性，与时间演化方程的抛物分类无关。全平面及有限穿孔平面都是上述直接证明的例子。下一节会说明，正基本核若在一个极点存在，就在所有极点存在；因此平面区域不会出现“某些极点属于 Green 情形、另一些属于抛物情形”的混合情况。

\ProseParagraphBreak
圆盘扩大的例子已经能看出发散机制。以原点为极点的经典核为 \(g_{B(0,R)}(z,0)=\log(R/|z|)\)，这个函数直接满足上一节的全部定义条件。固定 \(z\ne0\)，随 \(R\to\infty\) 核值趋于正无穷。每个有界圆盘都可容纳有限点源响应，但把零边界不断推向无穷远后，在固定观察点得不到有限极限。

\ProseParagraphBreak
维数在这里有实际作用。在三维，正基本解与 \(|x-w|^{-1}\) 成比例，它在整个空间保持正且在无穷远趋零；二维的基本解则是对数。因而“把边界推到无穷远”是否仍保留有限正核，与基本解的尺度有关，不能把平面的圆周平均论证原样移到高维。

\subsection{覆盖类型与 Green 存在性的区别}
\label{b3:subsec:170204}

Riemann 曲面论通常把万有覆盖为圆盘的曲面称为双曲型；有些势论文献也用“双曲型”表示存在正 Green 函数。两种说法描述不同条件，本章统一用“Green 域”指核的存在性。

\ProseParagraphBreak
去心圆盘是有界域，所以是 Green 域；它的万有覆盖也是圆盘。另一方面，\cref{b3:prop:12-planar-hyperbolic}说明，补集至少含两个点的平面域具有圆盘万有覆盖，故 \(\mathbb C\setminus\{0,1\}\) 属于这种覆盖类型；该结论沿用第12章已标明的一般统一化输入。但\cref{b3:prop:17-finite-punctures-parabolic}表明它没有正 Green 核。这个例子同时具有圆盘覆盖与势论抛物性。已知万有覆盖的共形类型，尚不足以判断点源在原域内是否产生有限正响应。

\ProseParagraphBreak
单连通情形则可以直接转移。若 \(D\subsetneq\mathbb C\) 为单连通域，Riemann 映照定理给出共形双射 \(F:D\to\mathbb D\)。固定 \(w\in D\)，取圆盘的经典 Green 核，并令
\[
 q(z)=g_{\mathbb D}(F(z),F(w)).
\]
它在极点外正且调和。因 \(F'(w)\ne0\)，在 \(w\) 附近有
\[
 F(z)-F(w)=(z-w)A(z),\quad A(w)=F'(w)\ne0,
\]
其中 \(A\) 全纯且局部无零。因此复合核加上 \(\log|z-w|\) 后，只新增调和项 \(-\log|A(z)|\)，仍具有所需奇性。它是 \(\mathcal Q_D(w)\) 的成员，故任意真单连通平面域都是 Green 域。这个存在性判断只用内部共形双射，不要求 \(F\) 连续延拓到通常拓扑边界。

\ProseParagraphBreak
若只有覆盖映射而没有双射，原域的一个极点会有多个原像。把原域的核拉回覆盖空间，便在每一个原像处产生对数奇性，不能把这个拉回误认为覆盖空间上只有一个极点的 Green 核。这正是多连通情形比共形双射多出的困难。

\subsection{与万有覆盖的联系}
\label{b3:subsec:170205}

设 \(\pi:\mathbb D\to D\) 为万有覆盖，固定 \(w\in D\) 及一个原像 \(\widetilde w\)。第11章的覆盖理论说明，纤维由覆盖变换群 \(\Gamma\) 的轨道 \(\{\gamma\widetilde w:\gamma\in\Gamma\}\) 给出；该集合离散且可数。于是一个自然的候选表达是
\[
 \sum_{\gamma\in\Gamma}
 g_{\mathbb D}(\widetilde z,\gamma\widetilde w).
\]
每项都非负；有限部分和只处理有限多个点源，而完整纤维要求控制无穷项。即使每个单项都有限，和仍可能发散。离散性只保证内部紧集包含有限个极点，并不控制趋近单位圆周的无穷多个极点的总贡献。

\ProseParagraphBreak
要用这一表达构造原域核，至少还须证明：部分和在极点轨道外局部一致收敛；在每个原像附近只留下一个系数为一的对数奇性；和对覆盖变换不变，从而能够下降。要进一步识别为最小核，还需最小性比较。这些问题各有作用，不能用“按群求和”一句话取代。本章的一般构造采用域内穷竭，因而不以该级数公式为前提。

\ProseParagraphBreak
从这里也可理解为何前述两种分类会分开。圆盘覆盖给出每个单独源点的正核，却没有自动给出整个轨道的求和控制；势论中的有限性正是额外要求。下一节把这种要求改写成更便于检验的条件：在一个固定非极点位置，穷竭核的数值有界。Harnack 链随后会把这个一点条件传播到所有内部紧集。

% !TeX root = ../../Book3.tex
\section{Green 函数的构造}
\label{b3:sec:1703}

先在规则有界域构造核，再逐步放大区域。单调性与 Harnack 原理控制极限，最小性使构造不依赖所选穷竭。

\subsection{穷竭}
\label{b3:subsec:170301}

任意平面区域可取 \(D_j\) 光滑有界连通子域，\(\overline{D_j}\subset D_{j+1}\)，且并为 \(D\)。例如先取到补集距离至少 \(1/j\)、模不超过 \(j\) 的紧部分，添加连接路径，再在内部用有限圆盘和光滑边界逼近。极点包含在首项之后。

若 \(D_j\subset D_{j+1}\)，核之差在极点可去，且在 \(\partial D_j\) 非负；最大值原理给 \(g_{D_j}\le g_{D_{j+1}}\)。这是真正的单调量。不能直接说边界值逐项收敛，因为边界本身在移动。

\subsection{Perron 方法}
\label{b3:subsec:170302}

对于光滑有界 \(D_j\)，用 Perron 法解边值 \(\log|z-w|\)，得调和函数 \(h_j\)，令 \(g_j=-\log|z-w|+h_j\)。小圆上它趋正无穷、外边界为零，最大值原理给严格正性。

\begin{theorem}[Green 核的穷竭构造]\label{b3:thm:17-green-exhaustion}
若 \(g_j(z_0,w)\) 对某个 \(z_0\ne w\) 有界，则 \(g_j\) 局部一致收敛于有限正 Green 核；该核是所有具有同一对数奇性的正调和核中的最小者，因而与穷竭无关。
\end{theorem}
\begin{proof}
极点外用 Harnack 收敛得到调和极限。在极点周围固定一小圆，边界上的极限有界；\(h_j=g_j+\log|z-w|\) 在圆盘内调和，最大值原理把边界界传到内部。Harnack 或 Poisson 表示使 \(h_j\) 在较小圆盘收敛到调和函数，所以奇性保持。
任取竞争正核 \(G\)，\(G-g_j\) 在极点可去，且在 \(\partial D_j\) 非负，因此 \(g_j\le G\)。取极限得最小性。不同穷竭产生的有限核彼此作为竞争核，故相等。若一个穷竭发散而另一个有限，后者又给前者共同上界，矛盾。
\end{proof}
有限核的对称性由光滑有界核的对称性取极限得到。

\subsection{圆盘与半平面}
\label{b3:subsec:170303}

圆盘与上半平面的显式核为
\[
g_{B(0,R)}(z,w)=\log\left|\frac{R^2-\overline wz}{R(z-w)}\right|,
\quad
g_{\mathfrak H}(z,w)=\log\left|\frac{z-\overline w}{z-w}\right|.
\]
分子零点在区域外，故余项调和。边界上分子和分母的模相等，且区域内两模之比大于一。圆盘由唯一性，半平面由共形变换和最小性确认。

极点在圆盘中心时核简化成 \(\log(R/|z|)\)。上半平面公式是负对数点源与反射源之差；反射点必须为 \(\overline w\)，不能给反射源保留同号。

\subsection{比较原理}
\label{b3:subsec:170304}

\begin{proposition}\label{b3:prop:17-domain-monotonicity}
若 Green 域 \(D\subset E\)，则 \(g_D(z,w)\le g_E(z,w)\)。
\end{proposition}
\begin{proof}
\(E\) 的核限制到 \(D\) 是一个正竞争核，具有正确奇性，最小性立即给不等式。
\end{proof}
将边界推远会减弱零势约束，核因而增大。对有界域，这个原则用包围圆盘给有限上界；它也使某些无界域可由半平面或其他已知 Green 域控制。


\ProseParagraphBreak

穷竭极限还保留极点参数的局部稳定性。在两个不相交的相对紧小圆盘 \(U,V\Subset D\) 上，光滑穷竭核在两个变量各自调和且对称。固定源点先用 Harnack 控制观察点，再利用对称性在源点方向使用同一估计，就可在 \(U\times V\) 上得到共同上界及导数界。极限因此在远离对角线处连续，并在每个变量中调和。仅仅知道逐点单调极限有限，还不能未经此控制就交换源点极限或把核当成连续的双变量函数。

\ProseParagraphBreak

在圆盘里，核的正性可由一个恒等式直接核对：
\[
|R^2-\overline wz|^2-R^2|z-w|^2=(R^2-|z|^2)(R^2-|w|^2)>0.
\]
两点均在圆盘内部时右端严格正，故模之比大于一。上半平面的相应恒等式是 \(|z-\overline w|^2-|z-w|^2=4\operatorname{Im}z\operatorname{Im}w>0\)。这些实代数算式避免凭图形猜测反射点的位置和源的符号；也说明两核在光滑边界分别降到零。

\ProseParagraphBreak

若 \(D\subsetneq E\) 且两域连通，区域单调性一般是严格的。在 \(D\) 内两核差调和、非负；如果在一个内部非极点位置为零，强最小值原理使差在 \(D\) 恒零。对光滑域，如果 \(\partial D\) 有一个点位于 \(E\) 内，\(E\) 的核在那里严格正而 \(D\) 的核取零，便产生矛盾。对于粗糙域，删去一个极性孔可能不改变核，所以不能仅由集合严格包含就断言严格不等式。严格性需要边界差异在势论中具有实际影响。

% !TeX root = ../../Book3.tex
\section{Green 势、能量与容量}
\label{b3:sec:1704}

点源的位置由一个参数表示，分布源则由测度表示。把 Green 核按源测度积分，可以同时保留对数奇性、区域几何和零边值约束。本节先说明这种叠加在何种意义下解方程，再从核的边界导数恢复调和测度，最后讨论在一个紧集上怎样分配单位质量才能使相互作用能量最小。

\subsection{Green 势}
\label{b3:subsec:170401}

若有限多个源点 \(w_j\) 的质量为 \(a_j>0\)，叠加势就是 \(\sum_j a_jg_D(z,w_j)\)。将离散质量换成连续或奇异的分布时，有限和自然换成积分。对角上约定 \(g_D(w,w)=+\infty\)，于是原子仍保留点源的无穷奇性。

\Needspace{6\baselineskip}
\begin{definition}\label{b3:def:17-green-potential}
设 \(D\subseteq\mathbb C\) 为 Green 域，\(\mu\) 为 \(D\) 上有限正 Borel 测度。函数
\[
 G_D\mu(z)\coloneqq\int_Dg_D(z,w)\,\mathrm d\mu(w),
 \quad z\in D,
\]
称为 \(\mu\) 的\term[green-shi]{Green 势}[Green potential]，积分按非负扩展积分理解，允许值为 \(+\infty\)。
\end{definition}
\symidx{\(G_D\mu\)：正测度 \(\mu\) 的 Green 势}

核的联合下半连续性保证可测性，所以定义不需要先假定势有限。要把它看成一个分布，还须证明局部可积；下面先处理质量支集留在域内紧集上的情形。

\Needspace{6\baselineskip}
\begin{proposition}[紧支测度的 Green 势]\label{b3:prop:17-green-potential-distribution}
设 \(D\subseteq\mathbb C\) 为 Green 域，有限正 Borel 测度 \(\mu\) 的支集包含于某个 \(K\Subset D\)。则 \(G_D\mu\in L^1_{\mathrm{loc}}(D)\)，并满足分布等式
\begin{equation}\label{b3:eq:17-green-source}
 -\Delta G_D\mu=2\pi\mu.
\end{equation}
具体而言，对每个 \(\varphi\in C_c^\infty(D)\)，有
\[
 \int_DG_D\mu(z)\,(-\Delta\varphi(z))\,\mathrm dA(z)
 =2\pi\int_D\varphi(w)\,\mathrm d\mu(w).
\]
\end{proposition}
\begin{proof}
固定观察点紧集 \(L\Subset D\)。由\cref{b3:prop:17-kernel-regularity}，对数余项在内部紧集上有共同界；远离对角线的核又连续。因此存在 \(C\ge0\)，使
\[
 0\le g_D(z,w)\le C+\log^+\frac1{|z-w|}
 \quad(z\in L,\ w\in K),
\]
其中 \(\log^+t=\max\{\log t,0\}\)。这个控制中的常数与源点 \(w\) 无关。由于
\[
 \int_L\log^+\frac1{|z-w|}\,\mathrm dA(z)
 \le2\pi\int_0^1-s\log s\,\mathrm ds=\frac\pi2,
\]
Tonelli 定理给出
\[
 \int_LG_D\mu\,\mathrm dA
 \le\mu(K)\left(C\operatorname{area}(L)+\frac\pi2\right)<\infty.
\]

对测试函数 \(\varphi\)，取 \(L=\operatorname{supp}\varphi\)，上述估计乘以 \(\|\Delta\varphi\|_\infty\) 便给出换序所需的绝对可积性。每个源点的核满足 \(-\Delta_zg_D(z,w)=2\pi\delta_w\)：这是\cref{b3:prop:15-log-fundamental}与核的调和余项给出的分布等式。因此
\[
 \begin{aligned}
 \int_DG_D\mu(z)(-\Delta\varphi(z))\,\mathrm dA(z)
 &=\int_K\left(\int_Dg_D(z,w)(-\Delta\varphi(z))\,\mathrm dA(z)\right)\mathrm d\mu(w)\\
 &=2\pi\int_K\varphi(w)\,\mathrm d\mu(w).
 \end{aligned}
\]
\end{proof}

局部可积只保证势几乎处处有限。例如 \(G_D\delta_w=g_D(\cdot,w)\) 在 \(w\) 处为正无穷，仍然满足上面的分布方程。对支集不紧的有限测度，本节保留势的扩展积分定义；有关局部可积性或无穷远收敛的进一步结论需要另外的控制。

\ProseParagraphBreak

光滑密度会把对数奇性平均掉。此时还可把分布解恢复成经典解，并证明它连续取零边值。

\Needspace{6\baselineskip}
\begin{proposition}[光滑内部源的零边值解]\label{b3:prop:17-green-poisson-source}
设 \(D\subseteq\mathbb C\) 为有界区域，其边界由有限条互不相交的正则实解析 Jordan 曲线组成，\(f\in C_c^\infty(D;\mathbb R)\)。则
\[
 u_f(z)\coloneqq\frac1{2\pi}\int_Dg_D(z,w)f(w)\,\mathrm dA(w)
\]
在 \(D\) 内光滑，连续延拓到 \(\overline D\)，并且是
\[
 -\Delta u=f\quad\text{于 }D,
 \quad u=0\quad\text{于 }\partial D
\]
在 \(C^2(D)\cap C(\overline D)\) 中的唯一解。函数 \(u_f\) 还在每段边界附近具有光滑调和延拓。
\end{proposition}
\begin{proof}
把 \(f\) 在域外延为零，记 \(K=\operatorname{supp}f\)。核的对数控制保证积分在每个点绝对收敛。写
\[
 g_D(z,w)=-\log|z-w|+h(z,w),
\]
其中 \(h(\cdot,w)\) 在填补 \(z=w\) 的可去点后于全域调和。于是 \(u_f=N_f+H_f\)，其中
\[
 \begin{aligned}
 N_f(z)&=-\frac1{2\pi}\int_{\mathbb C}\log|z-w|\cdot f(w)\,\mathrm dA(w),\\
 H_f(z)&=\frac1{2\pi}\int_Kh(z,w)f(w)\,\mathrm dA(w).
 \end{aligned}
\]
对 \(N_f\) 作换元 \(t=z-w\)，可将各阶导数放到 \(f(z-t)\) 上；当 \(z\) 留在固定紧集时，积分只涉及共同有界的 \(t\) 集合，而 \(\log|t|\) 局部可积。因此 \(N_f\) 光滑，且对数基本解给 \(-\Delta N_f=f\)。对 \(H_f\)，\cref{b3:prop:17-kernel-regularity}的共同余项界与\cref{b3:prop:15-interior-derivatives}给出关于 \(z\) 的各阶导数的共同局部界，故可在积分号下求导，得到 \(H_f\) 光滑且调和。这证明了内部方程。

还须控制 \(z\) 接近边界时的源点积分。选定 \(w_0\in D\)，用内部有限条路径及小圆盘把 \(K\) 与 \(w_0\) 连在一个相对紧的连通邻域中。在足够薄的边界邻域内，观察点 \(z\) 不落入这些圆盘。函数 \(w\mapsto g_D(z,w)\) 在这组圆盘上正且调和，所以\cref{b3:prop:15-harnack-chain}的有限链比较给出
\[
 \sup_{w\in K}g_D(z,w)\le C_Kg_D(z,w_0).
\]
所用圆盘固定，常数不随 \(z\) 改变。右端在 \(z\to\partial D\) 时趋于零，因而
\[
 |u_f(z)|\le\frac{C_K}{2\pi}\|f\|_{L^1(D)}g_D(z,w_0)\longrightarrow0.
\]
这给出连续零边值。源 \(f\) 的支集与边界有正距离，故 \(u_f\) 在边界内侧已经调和；\cref{b3:sec:1701}的解析边界拉直与奇反射使它光滑延拓到每段边界的另一侧。最后，两解之差调和并连续取零边值，最大值原理给出唯一性。
\end{proof}

式\eqref{b3:eq:17-green-source}解释了这里的系数：Green 势本身产生 \(2\pi f\)，方程 \(-\Delta u=f\) 的解须除以 \(2\pi\)。后面的边界表示采用同一个归一化。

\subsection{调和测度}
\label{b3:subsec:170402}

上一章从边界解算子的正性得到调和测度，没有要求边界存在法向量。对于解析边界，Green 核的反射正则性允许把这个抽象测度写成弧长密度。先对具有边界导数的函数证明积分公式，再以一致逼近处理一般连续边值。

\Needspace{6\baselineskip}
\begin{proposition}[Green 第二表示公式]\label{b3:prop:17-green-representation}
设 \(D\subseteq\mathbb C\) 为有界区域，其边界由有限条互不相交的正则实解析 Jordan 曲线组成；实函数 \(u\) 在 \(\overline D\) 的一个开邻域上属于 \(C^2\)。则对每个 \(z\in D\)，
\begin{equation}\label{b3:eq:17-green-representation}
 u(z)=-\frac1{2\pi}\int_{\partial D}u(\xi)\partial_{n_\xi}g_D(\xi,z)\,\mathrm ds(\xi)
 -\frac1{2\pi}\int_Dg_D(\xi,z)\Delta u(\xi)\,\mathrm dA(\xi),
\end{equation}
其中 \(n_\xi\) 是区域的外单位法向，\(\mathrm ds\) 为未归一化弧长测度。
\end{proposition}
\begin{proof}
固定 \(z\)，记 \(g(\xi)=g_D(\xi,z)\)，并从 \(D\) 删去闭圆盘 \(\overline{B(z,\varepsilon)}\)。核在外边界附近的光滑性由\cref{b3:sec:1701}的反射论证保证；在删圆后的闭域上，\(u\) 与 \(g\) 都有 Green 第二公式所需的正则性。由于 \(\Delta g=0\)，有
\[
 -\int_{D\setminus\overline{B(z,\varepsilon)}}g\Delta u\,\mathrm dA
 =\int_{\partial D}u\partial_ng\,\mathrm ds
 +\int_{|\xi-z|=\varepsilon}
       (u\partial_ng-g\partial_nu)\,\mathrm ds.
\]
外边界上 \(g=0\)。内圆的外法向是删圆余域的外法向，指向圆心，所以
\[
 g=-\log\varepsilon+O(1),
 \quad \partial_ng=\varepsilon^{-1}+O(1),
 \quad u(\xi)=u(z)+O(\varepsilon).
\]
内圆第一项的积分趋于 \(2\pi u(z)\)，第二项的绝对值为 \(O(\varepsilon|\log\varepsilon|)\)，趋于零。面积积分也收敛：极点附近的对数可积，远离极点的闭域部分上核有界。令 \(\varepsilon\downarrow0\) 并移项即得公式。
\end{proof}

内圆的法向决定公式中的符号。取 \(u\equiv1\)，面积项消失，边界项须给出一；取零边值函数，边界项消失，面积项则与上一小节的源表示相合。

\Needspace{6\baselineskip}
\begin{proposition}[解析边界上的调和测度密度]\label{b3:prop:17-harmonic-measure-density}
设 \(D\subseteq\mathbb C\) 为有界区域，其边界由有限条互不相交的正则实解析 Jordan 曲线组成。对 \(z\in D\)，定义
\[
 q_z(\xi)\coloneqq-\frac1{2\pi}\partial_{n_\xi}g_D(\xi,z)
 \quad(\xi\in\partial D).
\]
则 \(q_z\) 连续、非负，且 \(\int_{\partial D}q_z\,\mathrm ds=1\)。该域的调和测度满足
\begin{equation}\label{b3:eq:17-harmonic-measure-density}
 \mathrm d\omega_D^z(\xi)=q_z(\xi)\,\mathrm ds(\xi).
\end{equation}
因此，对任意实值连续边界数据 \(\varphi\)，其经典调和解为
\[
 H_\varphi(z)=\int_{\partial D}\varphi(\xi)q_z(\xi)\,\mathrm ds(\xi).
\]
\end{proposition}
\begin{proof}
反射使固定源点的核在边界附近光滑，故 \(q_z\) 连续。沿边界的内法向进入区域时核非负，而边界值为零，因而
\[
 \partial_{n_\xi}g_D(\xi,z)
 =\lim_{t\downarrow0}\frac{g_D(\xi,z)-g_D(\xi-tn_\xi,z)}{t}\le0.
\]
对常数函数使用\cref{b3:prop:17-green-representation}，便得 \(q_z\) 的积分为一。

每个解析边界点都满足局部外锥条件，所以上一章的正则性结果保证连续边值问题可解。为把表示公式用于 \(H_\varphi\)，对每条边界分量选一个实解析的正则周期参数 \(\gamma_j(t)\)，周期取 \(2\pi\)。连续函数 \(\varphi\circ\gamma_j\) 可由实三角多项式 \(p_{j,m}\) 一致逼近；这可用\cref{b3:subsec:160204}中已经证明的 Fej\'er 平均取得。由有限条分量组成的边值 \(\varphi_m(\gamma_j(t))=p_{j,m}(t)\) 因而在整个 \(\partial D\) 一致趋于 \(\varphi\)。记其调和解为 \(u_m\)。

这些特殊边值确实给出边界光滑性。固定一段解析弧，参数 \(\gamma_j\) 的局部全纯逆把它拉直为实轴。三角多项式 \(p_{j,m}(t)\) 是某个全纯函数 \(P_{j,m}(t)\) 在实轴上的实部：将其中的实系数正弦、余弦直接延拓到复参数即可。因此
\[
 v(\xi)=\operatorname{Re}P_{j,m}(\gamma_j^{-1}(\xi))
\]
在弧的邻域调和，并取到 \(u_m\) 的边值。差 \(u_m-v\) 连续取零弧值，奇反射使它调和延拓。这样 \(u_m\) 在每个边界点附近都有光滑延拓；有限覆盖与光滑分片也给出闭域邻域上的 \(C^2\) 延拓。对 \(u_m\) 使用 Green 表示，面积项因调和性消失，得到
\[
 u_m(z)=\int_{\partial D}\varphi_m(\xi)q_z(\xi)\,\mathrm ds(\xi).
\]

最大值原理给
\[
 \|u_m-H_\varphi\|_{C(\overline D)}
 \le\|\varphi_m-\varphi\|_{C(\partial D)}\longrightarrow0.
\]
另一方面，\(q_z\,\mathrm ds\) 是概率测度，故边界积分的误差也不超过同一个一致范数。取极限得到连续数据的表示，再由\cref{b3:thm:16-harmonic-measure}中表示测度的唯一性识别出 \(\omega_D^z\)。
\end{proof}

一般连续边值的调和解不由此获得边界 \(C^1\) 正则性；证明只在三角多项式逼近的解上使用边界导数，再通过最大值原理取极限。导数落在固定内点的 Green 核上，正是这种表示适用于连续数据的原因。

\ProseParagraphBreak

在圆盘 \(D=B(0,R)\) 中，对显式 Green 核沿 \(\xi=Re^{\mathrm it}\) 的径向求导，得到
\[
 q_z(Re^{\mathrm it})
 =\frac{R^2-|z|^2}{2\pi R\,|Re^{\mathrm it}-z|^2}.
\]
乘以弧长元 \(R\,\mathrm dt\)，便恢复圆盘的 Poisson 权重。这里的 \(R\) 因子来自弧长与角参数的区别。

若同时给定 \(f\in C_c^\infty(D;\mathbb R)\) 和 \(\varphi\in C(\partial D;\mathbb R)\)，则唯一的经典解由两项叠加给出：
\[
 u(z)=\int_{\partial D}\varphi(\xi)q_z(\xi)\,\mathrm ds(\xi)
 +\frac1{2\pi}\int_Dg_D(z,w)f(w)\,\mathrm dA(w).
\]
它满足 \(-\Delta u=f\) 和 \(u|_{\partial D}=\varphi\)。边界项分配边界数据，面积项叠加内部源。作为常数核验，圆盘上的函数 \(u(z)=(R^2-|z|^2)/4\) 满足 \(-\Delta u=1\) 且边值为零。它在闭域邻域为 \(C^2\)，所以直接由\cref{b3:prop:17-green-representation}得到
\[
 \int_{B(0,R)}g_D(z,w)\,\mathrm dA(w)
 =\frac\pi2(R^2-|z|^2).
\]
这次源 \(1\) 不具有内部紧支集，使用的是已经适用的 Green 第二表示。中心点还可直接计算 \(2\pi\int_0^R s\log(R/s)\,\mathrm ds=\pi R^2/2\)。

\subsection{平衡观点}
\label{b3:subsec:170403}

现在固定一个紧集 \(K\Subset D\)，允许单位质量在 \(K\) 上重新分配。势在每个源点的值再按源测度平均，便得到总相互作用能量。对角上的无穷值使过度集中付出无穷能量；极小问题因而会区分一个集合是否能容纳有限能量的质量分布。

\Needspace{6\baselineskip}
\begin{definition}\label{b3:def:17-green-energy}
设 \(D\subseteq\mathbb C\) 为 Green 域，\(K\Subset D\) 为非空紧集，\(\mathcal P(K)\) 为支持在 \(K\) 上的 Borel 概率测度全体。对任意支持在 \(K\) 上的有限正 Borel 测度 \(\mu\)，定义其\term[green-nengliang]{Green 能量}[Green energy]，并在概率测度上取下确界：
\[
 I_D(\mu)\coloneqq\iint_{K\times K}g_D(z,w)\,\mathrm d\mu(z)\,\mathrm d\mu(w),
 \quad V_D(K)\coloneqq\inf_{\mu\in\mathcal P(K)}I_D(\mu).
\]
积分允许为正无穷。若 \(V_D(K)<\infty\)，且 \(\mu_K\in\mathcal P(K)\) 满足 \(I_D(\mu_K)=V_D(K)\)，则称 \(\mu_K\) 为 \(K\) 相对于 \(D\) 的\term[pingheng-cedu]{平衡测度}[equilibrium measure]。
\end{definition}
\symidx{\(I_D(\mu)\)：测度 \(\mu\) 的 Green 能量}
\symidx{\(V_D(K)\)：紧集 \(K\) 上概率测度的最小 Green 能量}

\Needspace{6\baselineskip}
\begin{theorem}[平衡测度的存在与变分条件]\label{b3:thm:17-equilibrium-existence}
设 \(D\subseteq\mathbb C\) 为 Green 域，\(K\Subset D\) 为非空紧集。如果 \(K\) 上存在有限 Green 能量的概率测度，则平衡测度 \(\mu_K\) 存在。对每个有限能量概率测度 \(\nu\in\mathcal P(K)\)，有
\begin{equation}\label{b3:eq:17-equilibrium-variation}
 \int_KG_D\mu_K\,\mathrm d\nu\ge I_D(\mu_K)=V_D(K),
\end{equation}
左端允许为正无穷。
\end{theorem}
\begin{proof}
先说明概率测度为何能取到弱收敛子列。紧度量空间 \(K\) 上的 \(C(K;\mathbb R)\) 有可数稠密子集：对每个正整数 \(m\)，选一个有限 \(1/m\)-网 \(\{x_i\}\)，以
\[
 \psi_i(x)=\frac{(2/m-|x-x_i|)_+}{\sum_j(2/m-|x-x_j|)_+}
\]
得到连续分片。分母处处为正，\(\sum_i\psi_i=1\)，且 \(\psi_i(x)>0\) 时 \(|x-x_i|<2/m\)。函数 \(\sum_i f(x_i)\psi_i\) 因一致连续性一致趋于 \(f\)；再将有限个系数换成有理数，所有 \(m\) 的这些组合便给出可数稠密族 \(\{f_\ell\}\)。

给定 \(\mu_n\in\mathcal P(K)\)，各数列 \(\int f_\ell\,\mathrm d\mu_n\) 有界。逐个抽子列并取对角线后，它们都收敛。概率质量为一使一致逼近误差在积分下不增大，故对每个 \(f\in C(K;\mathbb R)\)，积分也收敛。其极限 \(\Lambda(f)\) 是正线性泛函，且 \(\Lambda(1)=1\)。第一卷定理11.4.1的 Riesz--Markov 正泛函表示适用于紧 Hausdorff 空间 \(K\)，给出一个概率测度 \(\mu\)，使 \(\Lambda(f)=\int f\,\mathrm d\mu\)。于是选出的子列弱收敛于 \(\mu\)。

接着处理核的奇性。令 \(g_M=\min\{g_D,M\}\)，其中 \(M>0\)，对角值取 \(M\)。由\cref{b3:prop:17-kernel-regularity}，核在对角外连续，而内部紧集上的共同余项界给出
\[
 g_D(z,w)\ge-\log|z-w|-C_K.
\]
因此靠近对角线时 \(g_M\) 恒等于 \(M\)，从而 \(g_M\in C(K\times K)\)。

若 \(\mu_n\) 弱收敛于 \(\mu\)，则它们的积测度也对连续测试函数收敛。为看清这里没有遗漏双变量测试，给定 \(F\in C(K\times K)\)，用上面的细分片构造
\[
 F_m(z,w)=\sum_{i,j}F(x_i,x_j)\psi_i(z)\psi_j(w).
\]
一致连续性给 \(\|F_m-F\|_\infty\to0\)，每项的积分则分解为两个弱收敛积分的乘积。先固定 \(m\) 令 \(n\to\infty\)，再令 \(m\to\infty\)，便得到积测度的收敛。特别地，
\[
 \iint g_M\,\mathrm d\mu\,\mathrm d\mu
 =\lim_n\iint g_M\,\mathrm d\mu_n\,\mathrm d\mu_n
 \le\liminf_n I_D(\mu_n).
\]
再让 \(M\uparrow\infty\)，单调收敛给 \(I_D(\mu)\le\liminf_n I_D(\mu_n)\)。对一列能量趋于 \(V_D(K)\) 的概率测度取上述子列，就得到达到有限极小值的 \(\mu_K\)。

最后记 \(V=I_D(\mu_K)\)、\(J=\int G_D\mu_K\,\mathrm d\nu\)。若 \(J=\infty\)，所需不等式已经成立。若 \(J<\infty\)，对 \(0<t<1\) 取合法概率测度 \(\mu_t=(1-t)\mu_K+t\nu\)。正核的 Tonelli 换序与对称性给
\[
 I_D(\mu_t)=(1-t)^2V+2t(1-t)J+t^2I_D(\nu)\ge V.
\]
减去 \(V\)、除以 \(2t\)，再令 \(t\downarrow0\)，得到 \(J\ge V\)。各量在这一步均有限，所以没有对无穷能量作形式求导。
\end{proof}

变分条件比较的是对竞争测度的平均势。它不直接断言势在紧集的每一点都相等；特别是单点质量可能具有无穷能量，不能任取 \(\nu=\delta_z\) 再把有限能量变分公式当作逐点等式。一般集合上的精细平衡势结论还涉及零容量例外集，本节只在下面的圆盘例子中求出处处可核验的势。

\subsection{容量的联系}
\label{b3:subsec:170404}

能量越小，单位质量越容易分散。因此用最小能量的倒数衡量紧集在给定外域中的容纳能力。这个量保留外域：相同的内部集合，离零势边界的距离不同，其 Green 能量也不同。

\Needspace{6\baselineskip}
\begin{definition}\label{b3:def:17-green-capacity}
设 \(D\subseteq\mathbb C\) 为 Green 域，\(K\Subset D\) 为非空紧集。其\term[green-rongliang]{Green 容量}[Green capacity]定义为
\[
 \operatorname{Cap}_D(K)\coloneqq\frac1{V_D(K)},
 \quad \frac1{+\infty}=0.
\]
\end{definition}
\symidx{\(\operatorname{Cap}_D(K)\)：紧集 \(K\) 相对于 Green 域 \(D\) 的 Green 容量}

这个倒数不会因分母为零而失去意义。正连续函数 \(\min\{g_D,1\}\) 在 \(K\times K\) 上有严格正的最小值 \(a_K\)，故所有概率测度的能量至少为 \(a_K\)。若 \(K\) 只是有限点集，每个概率测度都有原子，对角奇性使其能量为无穷，因而容量为零。集合增大使可选测度增多，所以容量不减；外域增大使核增大，最小能量不减，容量便不增。这两个方向与上一章的相对 Dirichlet 容量一致。

\ProseParagraphBreak

下面完整计算同心闭圆盘的容量。先建立所需的能量非负性；它不能仅从核逐点为正推出，因为参与比较的是符号测度。

\Needspace{6\baselineskip}
\begin{lemma}[圆盘上的符号 Green 能量]\label{b3:lem:17-disk-energy-positive}
设 \(D=B(0,R)\)、\(R>0\)，\(\eta\) 是支持在某个内部紧集上的有限实符号 Borel 测度，并且
\[
 \iint g_D(z,w)\,\mathrm d|\eta|(z)\,\mathrm d|\eta|(w)<\infty.
\]
则绝对收敛的符号能量满足
\[
 I_D(\eta)\coloneqq\iint g_D(z,w)\,\mathrm d\eta(z)\,\mathrm d\eta(w)\ge0.
\]
\end{lemma}
\begin{proof}
先设 \(\mathrm d\eta=f\,\mathrm dA\)，其中 \(f\in C_c^\infty(D;\mathbb R)\)。令 \(V=G_D\eta\)。\cref{b3:prop:17-green-poisson-source}给 \(-\Delta V=2\pi f\)、零边值以及闭域上的光滑性。分部积分的边界项因 \(V=0\) 消失，于是
\[
 I_D(\eta)=\int_DVf\,\mathrm dA
 =\frac1{2\pi}\int_DV(-\Delta V)\,\mathrm dA
 =\frac1{2\pi}\int_D|\nabla V|^2\,\mathrm dA\ge0.
\]

一般情形下，取非负径向函数 \(\rho\in C_c^\infty(B(0,1))\)，使 \(\int\rho\,\mathrm dA=1\)，并令 \(\rho_\varepsilon(a)=\varepsilon^{-2}\rho(a/\varepsilon)\)。把 \(\eta\) 延为平面上的测度，记其磨光为 \(\eta_\varepsilon\)。当 \(\varepsilon\) 充分小时，磨光支集仍在同一个内部紧集上，因而已知 \(I_D(\eta_\varepsilon)\ge0\)。换序把这个能量写成
\[
 I_D(\eta_\varepsilon)=\iint g_\varepsilon(z,w)\,\mathrm d\eta(z)\,\mathrm d\eta(w),
\]
其中
\[
 g_\varepsilon(z,w)
 =\iint g_D(z+a,w+b)\rho_\varepsilon(a)\rho_\varepsilon(b)
            \,\mathrm dA(a)\,\mathrm dA(b).
\]
Green 核在每个变量中超调和，包括对角上的正无穷值。径向非负磨光是圆周平均的非负加权平均，因此超平均不等式可逐变量使用，得到
\[
 0\le g_\varepsilon(z,w)\le g_D(z,w).
\]
所有平移点都留在 \(D\) 内，故这里不涉及把核延到域外。

由对角外的联合连续性，\(g_\varepsilon(z,w)\to g_D(z,w)\) 当 \(z\ne w\)。有限能量假设还排除了 \(|\eta|\) 的原子：若一点质量为 \(a>0\)，对角贡献就使能量至少为 \(a^2M\)，对任意 \(M>0\) 均成立。于是对角线在 \(|\eta|\otimes|\eta|\) 下零测，因为其测度等于 \(\int|\eta|(\{z\})\,\mathrm d|\eta|(z)=0\)。共同控制 \(g_D\) 对这个积测度可积，支配收敛便给
\[
 I_D(\eta)=\lim_{\varepsilon\downarrow0}I_D(\eta_\varepsilon)\ge0.
\]
同一可积控制也使前面的换序对符号测度合法。
\end{proof}

\Needspace{6\baselineskip}
\begin{proposition}[同心闭圆盘的平衡测度与容量]\label{b3:prop:17-disk-capacity}
设 \(0<r<R\)，\(D=B(0,R)\)、\(K=\overline{B(0,r)}\)。令 \(\sigma_r\) 为圆周 \(|w|=r\) 上的均匀概率测度，即 \(\mathrm d\sigma_r(re^{\mathrm it})=\mathrm dt/(2\pi)\)。则
\begin{equation}\label{b3:eq:17-circle-equilibrium-potential}
 G_D\sigma_r(z)=
 \begin{cases}
  \log(R/r),&|z|\le r,\\
  \log(R/|z|),&r<|z|<R.
 \end{cases}
\end{equation}
测度 \(\sigma_r\) 是 \(K\) 的平衡测度，并且
\[
 V_D(K)=\log(R/r),
 \quad \operatorname{Cap}_D(K)=\frac1{\log(R/r)}.
\]
\end{proposition}
\begin{proof}
圆盘核可写成
\[
 g_D(z,w)=\log R+\log\left|1-\frac{\overline wz}{R^2}\right|-\log|z-w|.
\]
若 \(|a|<1\)，由 \(\log(1-ae^{\mathrm it})\) 的幂级数取实部并平均，得到
\[
 \frac1{2\pi}\int_0^{2\pi}\log|1-ae^{\mathrm it}|\,\mathrm dt=0.
\]
这里级数绝对一致收敛，逐项积分有依据。分子项总满足 \(|a|=r|z|/R^2<1\)，故其平均为零。分母项在 \(|z|<r\) 时提出 \(r\)，在 \(|z|>r\) 时提出 \(|z|\)，同一公式分别给出平均 \(\log r\) 与 \(\log|z|\)。

在 \(|z|=r\) 上还须处理圆周积分的一个对数奇点。对 \(1/2\le s<1\)，有
\[
 |1-se^{\mathrm it}|^2=(1-s)^2+4s\sin^2(t/2)
 \ge2\sin^2(t/2),
 \quad |1-se^{\mathrm it}|\le2.
\]
因此其对数绝对值被 \(C+|\log|\sin(t/2)||\) 控制，后者在圆周可积。令 \(s\uparrow1\)，支配收敛将平均零公式传到 \(|a|=1\)。由此圆周上的分母平均也等于 \(\log r\)，证明了包括界面点在内的势公式。记 \(c=\log(R/r)\)，于是 \(I_D(\sigma_r)=c\)。

取任意 \(\nu\in\mathcal P(K)\)。若其能量无穷，下界 \(I_D(\nu)\ge c\) 已成立；以下设其能量有限。由于 \(G_D\sigma_r=c\) 在整个 \(K\) 成立，交叉能量为
\[
 \iint g_D(z,w)\,\mathrm d\nu(z)\,\mathrm d\sigma_r(w)=c.
\]
令 \(\eta=\nu-\sigma_r\)。不等式 \(|\eta|\le\nu+\sigma_r\) 与正核给
\[
 I_D(|\eta|)\le I_D(\nu)+2c+I_D(\sigma_r)<\infty.
\]
因此\cref{b3:lem:17-disk-energy-positive}适用，并给出
\[
 0\le I_D(\eta)
 =I_D(\nu)-2c+I_D(\sigma_r)=I_D(\nu)-c.
\]
下界在 \(\sigma_r\) 处达到，所以它是平衡测度，最小能量及容量公式随之成立。
\end{proof}

同一势也说明本节与上一章的归一化差异。令 \(v=G_D\sigma_r/\log(R/r)\)，则 \(v\) 在内圆盘恒为一，在圆环中为 \(\log(R/|z|)/\log(R/r)\)，且在外圆周为零。其 Dirichlet 能量为
\[
 \int_D|\nabla v|^2\,\mathrm dA
 =\frac{2\pi}{\log^2(R/r)}\int_r^R\frac{\mathrm ds}{s}
 =\frac{2\pi}{\log(R/r)}.
\]
这个分段函数本身不是光滑紧支试函数，\cref{b3:prop:16-annulus-capacity}已用截断和磨光把它逼近为合法竞争函数，并证明了匹配的下界。因此在这个例子中，
\[
 \operatorname{cap}(K,D)=2\pi\operatorname{Cap}_D(K).
\]
因子 \(2\pi\) 同时出现在点源方程、势的能量恒等式和两种容量之间，均由 \(-\Delta(-\log|z|)=2\pi\delta_0\) 决定。一般紧集上两种容量的对应还需要进一步的势论，本节的计算已经说明为什么不能只因名称相近就混用数值。

% !TeX root = ../../Book3.tex
\section{多连通区域}
\label{b3:sec:1705}

多连通域的调和函数仍可由边值确定，但共轭的周期多了有限个独立方向。边界分量的调和测度提供了调整这些周期的自然函数。

\subsection{周期}
\label{b3:subsec:170501}

调和函数 \(u\) 的周期为 \(\int_\gamma(-u_y\,\mathrm{d}x+u_x\,\mathrm{d}y)\)。闭形式使周期只依赖曲线同调类。若区域由 \(m+1\) 条互不相交实解析 Jordan 曲线围成，则绕 \(m\) 个内孔的曲线构成独立周期方向；任一闭曲线的周期为这些基本周期的整数线性组合。

\subsection{调和测度}
\label{b3:subsec:170502}

令 \(\omega_j\) 为第 \(j\) 个内边界取一、其余边界取零的调和解。其边界数据在各分量上连续，各边界弧可拉直并取局部线性障碍，所以 Perron 理论保证存在；减去该分量上的常数后作调和反射，还保证解在边界附近光滑。它们满足 \(0<\omega_j<1\)，连同外边界权重 \(\omega_0\) 有 \(\sum_{j=0}^m\omega_j=1\)。因此常数边界数据 \(c_j\) 的解为 \(\sum c_j\omega_j\)。

环域 \(r<|z|<R\) 的内边界权重是 \(\log(R/|z|)/\log(R/r)\)，其共轭周期为 \(-2\pi/\log(R/r)\)。这给出可计算的周期调整系数。

\subsection{调和共轭的障碍}
\label{b3:subsec:170503}

沿正向边界，\(\omega\) 的积分等于调和函数的外法向通量。令 \(A_{ij}=\int_D\nabla\omega_i\cdot\nabla\omega_j\,\mathrm{d}A\)。Green 第一公式把它等于第 \(i\) 个内边界上的 \(\omega_j\) 通量。

矩阵 \(A\) 对称正定：若 \(\sum c_j\omega_j\) 的梯度能量为零，则该组合恒常；外边界值为零使常数零，再看各内边界值知全部 \(c_j=0\)。因此基本周期可以由 \(\omega_j\) 的线性组合任意调整，拓扑障碍有明确的有限维形式。

\subsection{典范表示}
\label{b3:subsec:170504}

\begin{proposition}\label{b3:prop:17-period-decomposition}
在上述有限连通实解析边界有界域中，每个连续到边界的调和函数可写为
\[
u=\operatorname{Re}F+\sum_{j=1}^m c_j\omega_j,
\]
其中 \(F\) 单值全纯，\(c_j\) 唯一；\(F\) 只差纯虚常数。
\end{proposition}
\begin{proof}
用正定通量矩阵解线性方程，选择 \(c_j\) 消除 \(u\) 在全部基本曲线上的周期。剩余调和函数所有闭曲线周期为零，由\cref{b3:thm:15-global-conjugate}得到单值共轭和 \(F\)。若两组系数都可行，其差的周期向量为零，矩阵可逆使差为零；共轭唯一性给最后结论。
\end{proof}


\ProseParagraphBreak

基本周期的方向须统一。沿区域正向边界行进时，区域在左侧，外边界逆时针、内边界顺时针。若切向为 \((x',y')\)，外法向为 \((y',-x')\)，则
\[
(-u_y\,\mathrm{d}x+u_x\,\mathrm{d}y)(\gamma')=\nabla u\cdot n\,\mathrm{d}s.
\]
所以通量矩阵直接给的是顺时针内边界周期；若基本闭曲线统一选逆时针绕孔，则矩阵前须加负号。环域内边界权重的逆时针周期是 \(-2\pi/\log(R/r)\)，而其沿区域正向内边界通量为 \(2\pi/\log(R/r)\)，正好与正定能量相符。

\ProseParagraphBreak

有限维拓扑描述可以通过割线实现。从每个内边界取一条互不相交的简单弧连接到外边界，沿这些弧割开区域后得到单连通域。在割开的域中闭形式有原函数。沿第 \(j\) 条割线两侧的原函数差为绕该孔的周期，且因导数相同而为常数；所有这些差均零时，原函数才能连续拼回原区域。任意闭曲线跨越各条割线的有向次数为整数，其积分等于对应周期的整数和。这同时解释基本周期足够、为何有 \(m\) 个独立方向，以及周期清零后共轭的单值性。

\ProseParagraphBreak

典范表示中的 \(c_j\) 不是边界数据的简单平均。它由整个调和函数的周期决定。环域中取 \(u(z)=\log|z|\)，其周期为 \(2\pi\)，内边界权重周期为 \(-2\pi/\log(R/r)\)，故 \(c=-\log(R/r)\)，余项是常数 \(\log R\)。取 \(u(z)=\operatorname{Re}z\)，周期为零，所以 \(c=0\)，尽管两个边界上的函数值通常不相同。边界权重分解只适合各分量上常数的边值，而周期分解适合一般调和函数；两种线性分解的系数含义不同。

\ProseParagraphBreak

当边界孔数增加，矩阵的大小也增加；如果有无穷多个孔，当前有限维正定论证并未给出一个自动收敛的无穷级数表示。系数、能量和周期的求和需要额外估计。即使孔数有限，若边界粗糙到通常法向导数不存在，也应改用内部周期和弱能量方法，不能继续逐点积分边界法向导数。本节的实解析边界有界范围正是为了使拓扑方向和通量方向能够用同一个有限矩阵联系起来。

% 本章小结（不计入编号）
\begin{chaptersummary}
有界规则域的 Green 核由边值解构造，一般域则由单调穷竭识别其最小正核。正核不存在的平面反例说明无穷远约束不可省略，也说明覆盖共形类型不等于 Green 存在性。核叠加产生势，能量极小化产生容量。有限连通域的通量矩阵把共轭障碍化成可解的线性方程。
\end{chaptersummary}


\BookExercises

% 本题组自拟；各题新增训练动作见题意，未复制第三方原题。
\begin{exercise}\label{b3:ex:17-draft-01}
核验单位圆盘 \(g(z,0)\) 的 Riesz 测度符号与边界值。
\end{exercise}
\begin{solution}
\(g(z,0)=-\log|z|\)，所以 \(-\Delta g=2\pi\delta_0\)，边界零。作为超调和函数，其 \(\Delta g\) 非正；不能把它的源质量按次调和 Riesz 测度记成正。
\end{solution}

\begin{exercise}\label{b3:ex:17-draft-02}
求上半平面中 \(g(\mathrm{i}y,\mathrm{i}a)\)，\(a,y>0\)、\(a\ne y\)，并考察 \(y\to0\)。
\end{exercise}
\begin{solution}
为 \(\log((y+a)/|y-a|)\)，正，因为 \(y+a>|y-a|\)。当 \(y\to0\) 时比值趋一，核趋零；当 \(y\to a\) 时趋正无穷。
\end{solution}

\begin{exercise}\label{b3:ex:17-draft-03}
设 \(D\subset B(0,R)\) 且 \(0\in D\)。证明 \(g_D(z,0)\le\log(R/|z|)\)。
\end{exercise}
\begin{solution}
右侧为包围圆盘的正 Green 核，限制到 \(D\) 保持奇性，最小性给结论。上界同时证明 \(D\) 的穷竭核有限。
\end{solution}

\begin{exercise}\label{b3:ex:17-draft-04}
解释为何在上半平面给正 Green 核加 \(cy\)、\(c>0\) 后不再是最小 Green 核。
\end{exercise}
\begin{solution}
新增项调和、边界零、在极点无奇性，故局部条件仍成立；但原核逐点严格较小，因而新增核不满足最小性。该例反驳无界域仅以通常边界值刻画唯一性。
\end{solution}

\begin{exercise}\label{b3:ex:17-draft-05}
求环域两边界分别取 \(A,B\) 的径向调和解，并给其共轭周期。
\end{exercise}
\begin{solution}
解为 \(A+(B-A)\log(|z|/r)/\log(R/r)\)，周期为 \(2\pi(B-A)/\log(R/r)\)。只有 \(A=B\) 才有单值共轭。
\end{solution}

\begin{exercise}\label{b3:ex:17-draft-06}
若有限正测度 \(\mu\) 有一个原子，证明其 Green 能量无穷。
\end{exercise}
\begin{solution}
若 \(\mu(\{a\})=c>0\)，双积分在 \((a,a)\) 上有质量 \(c^2\)，而 \(g_D(a,a)=+\infty\)。用截断核 \(\min(g_D,N)\) 下界能量至少为 \(c^2N\)，令 \(N\to\infty\)。
\end{solution}

\begin{exercise}\label{b3:ex:17-draft-07}
证明 \(\sum_{j=0}^m\omega_j=1\)，并说明为何这不造成内边界通量矩阵奇异。
\end{exercise}
\begin{solution}
左侧在每个边界分量取一，唯一性使其等于一。内边界矩阵只含 \(j=1,\ldots,m\)，这些组合在外边界取零，故不能等于非零常数，正定性不受全体权重关系影响。
\end{solution}

\begin{exercise}\label{b3:ex:17-draft-08}
在单位圆盘中，令 \(\mu\) 为半径 \(r\) 圆周的均匀概率测度。计算 \(G\mu(0)\) 与 \(I(\mu)\)。
\end{exercise}
\begin{solution}
中心核为 \(-\log r\)，故 \(G\mu(0)=-\log r\)。显式核中分子的对数圆周平均为零，而 \(\log|re^{\mathrm{i}t}-re^{\mathrm{i}s}|\) 平均为 \(\log r\)，所以圆周上势同为 \(-\log r\)，能量也为该值；对角单点在弧长测度中零测， 对数奇性 可积。
\end{solution}
